% TOP_PROBLEM: {"id": 2978, "title": "Kirby Problem 4.102", "category_id": 7, "category": {"id": 7, "name": "topology", "display_name": "Topology", "description": "Properties preserved under continuous deformations.", "slug": "topology", "order_index": 7, "created_at": "2026-07-31T15:26:25.671Z"}, "status": "open"}
% FIRST_POSTED: null
% ENTRY_KIND: statement_only
% Imported from: batch13.tex; UnsolvedMath ID 2978
\documentclass[11pt]{article}
\usepackage[margin=25mm]{geometry}
\usepackage{fontspec}
\setmainfont{Latin Modern Roman}
\usepackage{amsmath,amssymb,mathtools}
\usepackage{unicode-math}
\setmathfont{Latin Modern Math}
\usepackage{xurl,hyperref}
\hypersetup{colorlinks=true,urlcolor=blue}
\setcounter{secnumdepth}{0}
\setlength{\parindent}{0pt}
\setlength{\parskip}{5pt}
\setlength{\emergencystretch}{3em}
\newcommand{\sref}[2]{\hyperlink{source:#1:#2}{[S#2]}}
\newcommand{\cataloguescope}[1]{\paragraph{Recorded target.}#1}
\begin{document}
% UnsolvedMath ID 2978; source catalogue code KP-4.102
\setcounter{section}{2977}
\section{Kirby Problem 4.102}\label{2978}
{\small\sffamily UnsolvedMath ID 2978 · Topology}
\subsection{Definitions and mathematical statement}

\paragraph{Shared notation.}$\mathbb N=\{1,2,\ldots\}$, and $\mathbb Z,\mathbb Q,\mathbb R,\mathbb C$ denote the integers, rationals, reals and complex numbers. Any other conventions are specified locally.


Use $\mathbb{CP}^2$ with its Fubini--Study symplectic form $\omega_{FS}$. A singular symplectic curve is a real two-dimensional surface, possibly with several irreducible components, that is embedded and symplectic away from finitely many isolated singular points. At each singular point there is a local symplectomorphism of pairs to a complex plane algebraic curve in an open subset of $(\mathbb C^2,\omega_{\mathrm{std}})$, as specified in Source S3; having only a topologically algebraic link is not the local-model definition. Its topological type is the ambient isotopy class of its link, obtained by intersecting a sufficiently small three-sphere around the point with the curve. A cusp is locally irreducible, having one local branch. A rational cuspidal curve is homeomorphic to $S^2$ and has only these cuspidal singularities.

An equisingular symplectic isotopy is a continuous smooth family of singular symplectic surfaces, smooth away from the moving singular points, in which the number of singularities and the topological link type at each one stay fixed. If components are reducible, label them continuously so their labels are preserved throughout the family. \sref{2978}{3} The genus of a component means the genus of its smooth normalization. A complex curve is the zero set of homogeneous complex polynomials in the usual complex projective plane.
\paragraph{Statement recorded in the supplied source.}
(a) Is every symplectic rational cuspidal curve in $(\mathbb{CP}^2,\omega_{FS})$ equisingularly symplectically isotopic to a complex curve? (b) More generally, characterize the types that admit a symplectic representative not equisingularly symplectically isotopic to any complex curve. Here a type specifies the genus of every irreducible component and all the singularity link types. The second part permits reducible curves and locally reducible complex singularities. The source also asks whether every pair of equisingular singular symplectic surfaces in $\mathbb{CP}^2$ is equisingularly symplectically isotopic. It reports a negative answer to this latter universal question from equisingular complex Zariski pairs; this is recorded historical context rather than an assertion that the universal pair-isotopy question remains unsettled. \sref{2978}{2}
\subsection{Short English statement}
Ask whether every rational cuspidal symplectic curve is equisingularly isotopic to a complex curve, classify types admitting exceptions, and record the universal pair-isotopy question together with its known complex Zariski-pair counterexamples.
\subsection{Sources}
\begin{description}
\item[\hypertarget{source:2978:1}{S1}] ULAM.AI and the UnsolvedMath Contributors, UnsolvedMath: A Curated Collection of Open Mathematics Problems, record 2978 (KP-4.102). CC BY 4.0. \url{https://huggingface.co/datasets/ulamai/UnsolvedMath}.
\item[\hypertarget{source:2978:2}{S2}] R. İnanç Baykur, Robion C. Kirby and Daniel Ruberman, K3—A New Problem List in Low-Dimensional Topology, Mathematical Surveys and Monographs 295 (2026), author version, Problem 4.102 and Remarks, PDF page 274. \url{https://math.berkeley.edu/sites/default/files/surv-295-ruberman-watermarked-author-pdf.pdf}.
\item[\hypertarget{source:2978:3}{S3}] Marco Golla and Laura Starkston, The symplectic isotopy problem for rational cuspidal curves, Compositio Mathematica 158 (2022), 1595--1682; arXiv:1907.06787v3, Definitions 2.5 and 2.10, pages 8--9. \url{https://arxiv.org/pdf/1907.06787v3}.
\end{description}
\label{end:2978}
\end{document}
